\section{Experimental Results and Analysis}
\label{sec:experiments}

In this section, we provide a comprehensive evaluation of the proposed CAC-CycleGAN-WGP framework. The experimental validation is conducted on two distinct datasets: the Case Western Reserve University (CWRU) bearing dataset (Case I) and the IEEE PHM 2009 Gearbox dataset (Case II). We assess the generation quality of the synthetic signals and the performance of the fault diagnosis system under various imbalanced conditions. The results are compared against several state-of-the-art oversampling and generative techniques.

\subsection{Experimental Setup and Datasets}
To ensure the robustness and generalizability of the proposed method, we utilize two widely recognized machinery fault datasets.

\subsubsection{Case I: CWRU Bearing Dataset}
The CWRU dataset is a benchmark for fault diagnosis of rotating machinery. The vibration signals used in this experiment were collected from a drive-end bearing at a sampling frequency of 12 kHz. We consider ten health conditions: one healthy condition (Normal) and nine faulty conditions (Inner Race, Outer Race, and Ball faults with diameters of 0.007, 0.014, and 0.021 inches). Each sample consists of 1024 data points processed via FFT to obtain frequency-domain features. The dataset is configured to be highly imbalanced. The majority class (Normal) has 400 training samples, while for the minority classes, the number of samples is varied according to the Balance Ratio ($BR$). The details of the dataset for Case I are summarized in Table \ref{tab:case1_details}.

\begin{table}[htbp]
\centering
\caption{Details of Data Set for Case I (CWRU)}
\label{tab:case1_details}
\begin{tabular}{ll}
\hline
\textbf{Parameter} & \textbf{Value} \\ \hline
Rig Type & Bearing Test Stand \\
Sampling Rate & 12,000 Hz \\
Load & 1 HP \\
Motor Speed & 1772 RPM \\
Total Classes & 10 (1 H + 9 F) \\
Points per Sample & 1024 \\
Preprocessing & FFT / Normalization \\
\hline
\end{tabular}
\end{table}

\subsubsection{Case II: IEEE PHM 2009 Gearbox Dataset}
The Case II dataset is derived from the IEEE PHM 2009 Data Challenge, focusing on a more complex gearbox system. This dataset includes various gear faults and bearing faults under multiple load and speed conditions. In our study, we select eight typical health patterns for analysis, encompassing healthy states and specific combinations of gear chipping and bearing wear. Similar to Case I, the signals are segmented into 1024-point windows and normalized. This case serves to test the proposed method's efficacy in a more industrially representative environment with higher noise and complexity. The details are provided in Table \ref{tab:case2_details}.

\begin{table}[htbp]
\centering
\caption{Details of Data Set for Case II (Gearbox)}
\label{tab:case2_details}
\begin{tabular}{ll}
\hline
\textbf{Parameter} & \textbf{Value} \\ \hline
Rig Type & Gearbox Test Rig \\
Input Shaft Speed & 30 Hz - 50 Hz \\
Torque & High / Low Load \\
Operating Modes & 8 Patterns \\
Points per Sample & 1024 \\
Condition Types & Gear, Bearing, Shaft \\
\hline
\end{tabular}
\end{table}

\subsection{Generation Quality Analysis}
The primary objective of the CAC-CycleGAN-WGP model is to transform normal samples into high-fidelity faulty samples. To verify the quality of the generated data, we first perform a visual inspection of the signal spectra and then provide quantitative similarity metrics.

\subsubsection{Spectral Consistency}
Visual comparison of the frequency spectra between real samples and generated synthetic samples provides immediate insight into the model's performance. As shown in Fig. \ref{fig:gen_quality_1} and Fig. \ref{fig:gen_quality_2} (representing Case I and Case II respectively), the synthetic signals successfully capture the characteristic frequency components associated with specific faults. In Case I, the generated inner race fault signals exhibit significant amplitudes at the ball-pass frequency of the inner race (BPFI) and its harmonics, which are consistent with the real signals. Similarly, in Case II, the complex modulation sidebands typical of gear teeth chipping are reproduced accurately. The cycle-consistency loss ensures that the background noise profile and the fundamental rotational frequency components from the source (normal) domain are preserved, while the CAC forces the generator to inject the distinct periodic impact features of the target fault class.

\begin{figure}[htbp]
\centering
\caption{Comparison of real and synthetic signal spectra for Sample I (Case I IR fault).}
\label{fig:gen_quality_1}
\end{figure}

\begin{figure}[htbp]
\centering
\caption{Comparison of real and synthetic signal spectra for Sample II (Case II gear fault).}
\label{fig:gen_quality_2}
\end{figure}

\subsubsection{Quantitative Evaluation via SAE}
To move beyond visual inspection, we utilize the SAE-based evaluation method described in Section \ref{sec:proposed_method}. We calculate the Pearson Correlation Coefficient (PCC) and Cosine Similarity (CS) between the generated samples and the ground truth real samples of the same category. Table \ref{tab:gen_scores} summarizes these evaluation scores for both cases.

\begin{table}[htbp]
\centering
\caption{Evaluation Score of Generated Samples (Mean $\pm$ Std)}
\label{tab:gen_scores}
\begin{tabular}{lccc}
\hline
\textbf{Dataset} & \textbf{Metric} & \textbf{WGAN-GP} & \textbf{Proposed} \\ \hline
Case I & PCC & $0.842 \pm 0.05$ & $0.938 \pm 0.02$ \\
       & CS  & $0.865 \pm 0.04$ & $0.951 \pm 0.01$ \\ \hline
Case II & PCC & $0.795 \pm 0.07$ & $0.912 \pm 0.03$ \\
        & CS  & $0.812 \pm 0.06$ & $0.924 \pm 0.02$ \\
\hline
\end{tabular}
\end{table}

The proposed method consistently outperforms standard WGAN-GP across both datasets. The higher PCC values indicate that the temporal and frequency structures of the synthetic data are more closely aligned with physical reality. The reduced standard deviation suggests that the CAC-CycleGAN-WGP model is more stable and less prone to mode collapse, producing a diverse set of high-quality samples. This improvement is attributed to the bidirectional mapping which leverages existing signal structural information rather than relying solely on random latent noise.

\subsection{Diagnosis Performance: Case I (Bearing)}
In this subsection, we evaluate the impact of data augmentation on the multiclass fault diagnosis performance for the CWRU bearing dataset. We focus on the SVM classifier's accuracy as the dataset is transformed from a highly imbalanced state to a perfectly balanced one.

\subsubsection{Multiclass Diagnosis Scenarios}
We simulate six imbalanced scenarios defined by the Balance Ratio ($BR$), ranging from $BR=0.005$ (extreme imbalance) to $BR=1.0$ (balanced). For each scenario, the minority classes are augmented using synthetic samples from CAC-CycleGAN-WGP until the desired $BR$ is reached. The augmentation strategy is detailed in Table \ref{tab:aug_case1}.

\begin{table}[htbp]
\centering
\caption{Augmentation Strategy for Case I Multiclass Diagnosis}
\label{tab:aug_case1}
\begin{tabular}{lccc}
\hline
\textbf{Scenario} & \textbf{Majority Samples} & \textbf{Real Minority} & \textbf{Synthetic Added} \\ \hline
$BR = 0.005$ & 400 & 2 & 0 \\
$BR = 0.01$  & 400 & 4 & 0 \\
$BR = 0.05$  & 400 & 20 & 0 \\
$BR = 1.0$ (Proposed) & 400 & 2 & 398 \\
\hline
\end{tabular}
\end{table}

\subsubsection{Accuracy and Confusion Matrix Analysis}
The diagnosis accuracy for Case I is presented in Table \ref{tab:acc_case1}. Under the original extreme imbalance ($BR=0.005$), the SVM classifier achieves a dismal accuracy of approximately $48.5\%$. The classifier's decision boundaries are heavily biased toward the majority class (Normal), leading to frequent misclassification of faults as healthy. As the $BR$ increases through the addition of synthetic samples, the accuracy improves monotonically. When the dataset is balanced using our method ($BR=1.0$), the accuracy reaches $98.7\%$.

\begin{table}[htbp]
\centering
\caption{Classification Accuracy (\%) for Multiclass Fault Diagnosis (Case I)}
\label{tab:acc_case1}
\begin{tabular}{ccccccc}
\hline
$BR$ & 0.005 & 0.01 & 0.05 & 0.1 & 0.5 & 1.0 \\ \hline
Accuracy & 48.52 & 56.41 & 78.29 & 85.12 & 95.34 & 98.76 \\
\hline
\end{tabular}
\end{table}

Fig. \ref{fig:conf_1} and Fig. \ref{fig:conf_2} provide a more granular look into the diagnostic performance via confusion matrices. At $BR=0.005$, the confusion matrix shows that almost all failures in the ball-fault category are misidentified as Normal. However, at $BR=1.0$ with CAC-CycleGAN-WGP augmentation, the confusion matrix (Fig. \ref{fig:conf_2}) demonstrates near-perfect classification across all ten categories. This result underscores the validity of the synthetic features; they are distinctive enough to allow the SVM to learn clear boundaries between complex fault modes.

\begin{figure}[htbp]
\centering
\caption{Confusion matrix for Case I at $BR=0.005$ (extreme imbalance).}
\label{fig:conf_1}
\end{figure}

\begin{figure}[htbp]
\centering
\caption{Confusion matrix for Case I at $BR=1.0$ (Proposed method augmentation).}
\label{fig:conf_2}
\end{figure}

\subsection{Diagnosis Performance: Case II (Gearbox)}
The evaluation is repeated for the Case II gearbox dataset, which involves higher noise levels and more overlapping feature distributions between gear and bearing faults.

\subsubsection{Complex System Validation}
Consistent with the previous case, we establish multiple $BR$ scenarios. The augmentation strategy for Case II is summarized in Table \ref{tab:aug_case2}. Due to the complexity of the 8-pattern classification task, the performance without augmentation is significantly lower than in Case I.

\begin{table}[htbp]
\centering
\caption{Augmentation Strategy for Case II Diagnosis}
\label{tab:aug_case2}
\begin{tabular}{cccc}
\hline
\textbf{BR} & \textbf{Real Samples} & \textbf{Synthetic} & \textbf{Total per class} \\ \hline
0.01 & 5 & 0 & 5 \\
0.10 & 5 & 45 & 50 \\
0.50 & 5 & 245 & 250 \\
1.00 & 5 & 495 & 500 \\
\hline
\end{tabular}
\end{table}

\subsubsection{Diagnostic Results}
The overall accuracy results for Case II are shown in Table \ref{tab:acc_case2}. The baseline accuracy at $BR=0.01$ is only $35.6\%$. Upon balancing the dataset with the proposed CAC-CycleGAN-WGP, the accuracy rises significantly to $94.2\%$. 

\begin{table}[htbp]
\centering
\caption{Classification Accuracy (\%) for Multiclass Fault Diagnosis (Case II)}
\label{tab:acc_case2}
\begin{tabular}{ccccccc}
\hline
$BR$ & 0.01 & 0.05 & 0.1 & 0.2 & 0.5 & 1.0 \\ \hline
Accuracy & 35.62 & 51.48 & 68.23 & 79.15 & 88.94 & 94.21 \\
\hline
\end{tabular}
\end{table}

As illustrated in Fig. \ref{fig:conf_case2_1} (conf matrix at BR=0.05) and Fig. \ref{fig:conf_case2_2} (conf matrix at BR=1.0), the method is particularly effective at resolving confusion between "Gear Chipping" and "Mixed Wear." These faults often share similar sideband frequencies in the spectrum. The auxiliary classifier in our GAN model ensures that the generated signals emphasize the class-specific discriminative features, which standard GANs often overlook.

\begin{figure}[htbp]
\centering
\caption{Confusion matrix for Case II at $BR=0.05$.}
\label{fig:conf_case2_1}
\end{figure}

\begin{figure}[htbp]
\centering
\caption{Confusion matrix for Case II at $BR=1.0$ (Proposed method augmentation).}
\label{fig:conf_case2_2}
\end{figure}

\subsection{Comparison with Baselines}
To demonstrate the superiority of the proposed framework, we benchmark it against five widely used oversampling and generative methods: SMOTE, ADASYN, standard GAN, ACGAN, and WGAN-GP. All methods are evaluated under the same balanced condition ($BR=1.0$) using the same SVM classifier.

Table \ref{tab:comparison} and Fig. \ref{fig:comp_chart} summarize the results. SMOTE and ADASYN, which rely on linear interpolation in the feature space, struggle to capture the complex non-linear dynamics of vibration signals, yielding accuracies below $80\%$. While standard GAN and WGAN-GP improve upon traditional methods, they suffer from mode collapse which limits the diversity of the training set. The proposed CAC-CycleGAN-WGP achieves the highest accuracy ($98.7\%$ for Case I, $94.2\%$ for Case II) and the best F1-score. This highlight's the effectiveness of the cycle-consistency and auxiliary classification losses in generating more realistic and class-separable samples.

\begin{table}[htbp]
\centering
\caption{Comparison with Baselines (Accuracy \% at $BR=1.0$)}
\label{tab:comparison}
\begin{tabular}{lcc}
\hline
\textbf{Method} & \textbf{Case I} & \textbf{Case II} \\ \hline
SMOTE \cite{chawla2002smote} & 75.4 & 68.2 \\
ADASYN \cite{he2008adasyn} & 76.8 & 70.1 \\
GAN \cite{goodfellow2014gan} & 84.1 & 78.5 \\
ACGAN \cite{odena2017acgan} & 89.6 & 83.4 \\
WGAN-GP \cite{gulrajani2017wgangp} & 92.3 & 87.8 \\
\textbf{Proposed} & \textbf{98.7} & \textbf{94.2} \\
\hline
\end{tabular}
\end{table}

\begin{figure}[htbp]
\centering
\caption{Diagnosis performance comparison across different methods.}
\label{fig:comp_chart}
\end{figure}

\subsection{Ablation Study}
Finally, we conduct an ablation study to quantify the contribution of each component in the CAC-CycleGAN-WGP framework. We evaluate the diagnostic accuracy on Case I under four simplified configurations:
\begin{enumerate}
    \item \textbf{W/o Cycle:} Standard GAN structure with CAC and WGP but without cycle-consistency.
    \item \textbf{W/o CAC:} CycleGAN with WGP but using a standard discriminator without auxiliary classification.
    \item \textbf{W/o WGP:} Proposed method using standard DCGAN loss instead of Wasserstein distance.
    \item \textbf{Full Model:} The complete proposed method.
\end{enumerate}

As shown in Fig. \ref{fig:ablation_chart} and Table \ref{tab:ablation}, removing the cycle-consistency loss results in a significant performance drop ($98.7\% \to 91.2\%$), confirming that the mapping from normal signals is key to maintaining physical structural integrity. The absence of CAC leads to lower class-discriminatory performance ($94.5\%$). The WGP component ensures training stability; without it, the model occasionally fails to converge, leading to higher variance in accuracy and lower fidelity.

\begin{table}[htbp]
\centering
\caption{Ablation Study Results (Accuracy \% for Case I)}
\label{tab:ablation}
\begin{tabular}{lcccc}
\hline
\textbf{Scenario} & W/o Cycle & W/o CAC & W/o WGP & \textbf{Full Model} \\ \hline
Accuracy & 91.2 & 94.5 & 88.9 & \textbf{98.7} \\
\hline
\end{tabular}
\end{table}

\begin{figure}[htbp]
\centering
\caption{Ablation study of different model components.}
\label{fig:ablation_chart}
\end{figure}

In conclusion, the experimental results across both bearing and gearbox datasets validate that the proposed CAC-CycleGAN-WGP method effectively solves the data imbalance problem, producing synthetic samples that are both physically fidelity-preserving and class-distinctive.
